\chapter{Исследовательская часть}

% вар.21 => n=5 

\section{Графовые модели для рекурсивного алгоритма}

На рисунках 4.1 -- 4.4 представлены графовые модели реализации рекурсивного алгоритма.

\begin{figure}[!htbp]
	\centering
	\includesvg[width=0.2\textwidth]{r_gu}
	\caption{Граф управления.}
\end{figure}

\begin{figure}[!htbp]
	\centering
	\includesvg[width=0.3\textwidth]{r_ig}
	\caption{Информационный граф.}
\end{figure}

\begin{figure}[!htbp]
	\centering
	\includesvg[width=0.7\textwidth]{r_oi}
	\caption{Операционная история для n = 5.}
\end{figure}

\begin{figure}[!htbp]
	\centering
	\includesvg[width=0.7\textwidth]{r_ii}
	\caption{Информационная история для n = 5.}
\end{figure}

\FloatBarrier

\section{Графовые модели для нерекурсивного алгоритма}

На рисунках 4.5 -- 4.8 представлены графовые модели реализации нерекурсивного алгоритма.

\begin{figure}[!htbp]
	\centering
	\includesvg[width=0.15\textwidth]{nr_gu}
	\caption{Граф управления.}
\end{figure}

\begin{figure}[!htbp]
	\centering
	\includesvg[width=0.15\textwidth]{nr_ig}
	\caption{Информационный граф.}
\end{figure}

\begin{figure}[!htbp]
	\centering
	\includesvg[width=0.7\textwidth]{nr_oi}
	\caption{Операционная история для n = 5.}
\end{figure}

\begin{figure}[!htbp]
	\centering
	\includesvg[width=0.7\textwidth]{nr_ii}
	\caption{Информационная история для n = 5.}
\end{figure}

\FloatBarrier

\section*{Вывод}

В данном разделе были построены графовые модели для рекурсивного и нерекурсивного алгоритмов.

\clearpage
